% Section 8.3: The Cherenkov Telescope Array Observatory
% Source: 8.3_ctao.md (consolidated, post-review)

\section{The Cherenkov Telescope Array Observatory}
\label{sec:8.3}

The cross-correlation technique described in the preceding section reaches its full potential when paired with an instrument whose energy coverage, sky coverage, and angular resolution are matched to the physical scales of the DM signal. The Cherenkov Telescope Array Observatory (CTAO) satisfies all three requirements. Its energy range, extending from 20~GeV to 300~TeV, spans precisely the regime where heavy WIMP candidates produce their signal; its planned extragalactic survey covers a quarter of the sky with approximately uniform exposure; and its angular resolution at the relevant energies exceeds that of Fermi-LAT by an order of magnitude. The following subsections introduce CTAO's key characteristics and its extragalactic survey strategy before connecting them to the sensitivity forecast that constitutes the paper in this chapter.

\subsection{From Fermi-LAT to CTAO}
\label{sec:8.3.1}

The analyses in Chapters~\ref{ch:04}--\ref{ch:07} were built on data from the Fermi Large Area Telescope (Fermi-LAT), which has operated as the premier instrument for indirect dark matter searches in the GeV band for over 15 years (see Section~\ref{sec:2.3}). As a space-based pair-conversion telescope, the LAT detects gamma rays from 20~MeV to beyond 300~GeV with a field of view of approximately 2.4~sr, surveying the entire sky every two orbits ($\sim 3$~hours)~\cite{Fermi-LAT:2009ihh}. Its effective area, however, is constrained by spacecraft payload limitations to roughly 1~m$^2$ at GeV energies, a bound common to all satellite-borne instruments. This ceiling in collecting area, combined with the steeply falling flux of gamma rays at high energies, makes Fermi-LAT poorly suited for probing WIMP candidates with masses well above 1~TeV.

The Cherenkov Telescope Array Observatory (CTAO) is designed to overcome exactly this limitation~\cite{CTAConsortium:2017dvg}. Rather than detecting photons directly in space, ground-based imaging atmospheric Cherenkov telescopes (IACTs) reconstruct very-high-energy gamma rays from the brief flashes of Cherenkov light produced when an incident photon strikes the upper atmosphere and initiates an electromagnetic shower. By sampling this shower at ground level, CTAO effectively uses the atmosphere as a calorimeter of continental scale, achieving effective areas orders of magnitude larger than any space-based detector. The observatory will consist of two arrays, one in the northern hemisphere and one in the south, built from three telescope types tailored to different energy regimes: the Large-Sized Telescopes (LSTs, sensitive at lower energies, field of view $4.3^\circ$), the Medium-Sized Telescopes (MSTs, field of view $7.5^\circ$--$7.7^\circ$), and the Small-Sized Telescopes (SSTs, fields of view $8.8^\circ$, deployed only in the south where Galactic observations are prioritized). Across the three telescope classes, CTAO covers a broad energy range from 20~GeV up to 300~TeV, with an angular resolution and sensitivity that improve by an order of magnitude over current-generation IACTs such as H.E.S.S., MAGIC, and VERITAS~\cite{CTAConsortium:2017dvg}.

The extension to TeV energies opens a qualitatively new window for dark matter searches. Heavy WIMP candidates with masses $m_\chi \gtrsim 1$~TeV produce gamma rays predominantly in the TeV band, a range where Fermi-LAT has negligible sensitivity. More importantly for cross-correlation studies, EBL absorption becomes significant above $\sim 1$~TeV, attenuating photons from sources at $z \gtrsim 0.1$--$0.2$ and effectively cutting off the astrophysical backgrounds. The LST and MST telescopes, built in both hemispheres, are optimised for extragalactic observations, where the unresolved emission receives its dominant DM contribution from structures at $z \lesssim 0.1$~\cite{Pinetti:2025hgd}. At energies above 20~GeV and increasingly toward the TeV regime, star-forming galaxies and misaligned AGN contribute negligibly to the unresolved background; only blazars, both High-Synchrotron-Peaked (HSP) and Low-to-Intermediate-Synchrotron-Peaked (LISP), remain as the principal astrophysical foreground. Their window functions peak at $z \sim 0.3$--$0.4$ at 50~GeV and are pushed down to $z \sim 0.1$--$0.2$ above 1~TeV by EBL attenuation --- a range that still sits above the DM-dominated window, which we exploit through the choice of the 2MASS/2MRS galaxy catalogs (Section~\ref{sec:8.2.3}). This residual separation between the DM and blazar window functions is the physical lever arm that makes the cross-correlation technique viable at TeV energies with CTAO.

\subsection{The Extragalactic Survey}
\label{sec:8.3.2}

Cross-correlation and anisotropy analyses require large sky maps with approximately uniform exposure. A pointed instrument with irregular coverage would introduce instrumental variations that could mimic or obscure genuine cosmological angular correlations. For CTAO, this requirement is met by the Extragalactic Survey (EGAL), one of the observatory's designated Key Science Projects~\cite{CTAConsortium:2017dvg,Pinetti:2025hgd}.

The EGAL survey targets a quarter of the sky, covering the region defined by $-90^\circ < l < 90^\circ$ and $|b| > 5^\circ$ in Galactic coordinates. The survey is planned over three years of operations, split between the two array sites: the Southern array will observe 15\% of the sky with 400 hours, while the Northern array will cover the remaining 10\% with 600 hours, for a total observation time of approximately 1000 hours~\cite{CTAConsortium:2017dvg}. The default observing strategy is a grid of telescope pointings spaced by approximately $3^\circ$ across the survey area. Individual pointings are observed for 0.51 hours in the south and 1.11 hours in the north, yielding an effective average exposure of approximately 3 hours per point within the survey boundary~\cite{Pinetti:2025hgd}. The wide fields of view of the individual telescopes --- up to $8.8^\circ$ for the SSTs and $7.5^\circ$--$7.7^\circ$ for the MSTs --- create deliberate overlaps between adjacent pointings. These overlaps smooth out the exposure variations that inevitably arise when stitching together individual observations, keeping the relative fluctuations in the exposure map to approximately 10\% within the survey region~\cite{Pinetti:2025hgd}.

A second observation scenario considered in the paper appears from the accumulated off-source data collected through years of CTAO pointed observations across the extragalactic sky. In the most optimistic projection, these off-source data sum to an effective exposure of approximately 50 hours per sky location --- roughly 25 times larger than the nominal EGAL exposure --- yielding a factor of $\sim 4$ improvement in sensitivity compared to the 3-hour scenario~\cite{Pinetti:2025hgd}. This scenario does not require a dedicated survey and could in principle be reached earlier in the observatory's operational lifetime, making it a practically relevant complement to the EGAL Key Science Project.

Uniform sky coverage is a technical requirement, not just an observational convenience, for the cross-correlation formalism. The angular power spectrum $C_\ell$ is computed from a masked, pixelised sky map, and anisotropies in the instrumental exposure cannot be distinguished from intrinsic anisotropies in the gamma-ray background without careful masking and exposure correction. At the multipoles $\ell \sim 10$--$100$ relevant for the DM cross-correlation signal, any residual exposure gradient couples into the signal through the survey window function and biases the recovered power spectra. The EGAL strategy's uniform grid, combined with the broad telescope fields of view, is designed to minimise this systematic effect, providing the map quality required for the sensitivity forecast presented in the following paper.

\subsection{Transition to the Paper}
\label{sec:8.3.3}

With the CTAO observatory and its extragalactic survey established as the observational platform, and the cross-correlation formalism of Section~\ref{sec:8.2} in hand, we are now in a position to ask how well CTAO can constrain dark matter through its correlation with the large-scale distribution of matter. The paper that follows presents a sensitivity forecast addressing this question directly. Using the angular cross-power spectrum between the CTAO extragalactic survey and the 2MASS galaxy catalog --- the low-redshift gravitational tracer identified in Section~\ref{sec:8.2.3} as optimal for DM searches --- we derive the expected sensitivity to both DM annihilation and decay across a wide range of particle masses, for the $b\bar{b}$, $W^+W^-$, $\tau^+\tau^-$, and $e^+e^-$ channels~\cite{Pinetti:2025hgd}.

The main results are twofold. For the astrophysical signal from unresolved blazars, the forecast shows that a purely astrophysical detection through cross-correlation with 2MASS remains marginal even after 50 hours of CTAO data --- SNR $\simeq 0.16$ for HSP blazars, compared to SNR $\simeq 0.9$ in auto-correlation (at the nominal 50-hour point-source detection threshold)~\cite{Pinetti:2025hgd}. This is a feature of the cross-correlation technique, not a limitation: by suppressing high-redshift blazar contributions, the method trades astrophysical sensitivity for background rejection, directly benefiting DM searches. For dark matter, cross-correlation with 2MASS reaches upper bounds on the annihilation cross-section of $\langle\sigma v\rangle \sim 10^{-23}$~cm$^3$~s$^{-1}$ and lower bounds on the decay lifetime of $\tau \sim 10^{27}$~s, for 50 hours of observation. These sensitivities improve by approximately a factor of 5 over the auto-correlation results at the same exposure, as shown in Figure~\ref{fig:bounds1}.

The cross-correlation limits are competitive with constraints from targeted pointed observations --- dwarf spheroidal galaxies for annihilation and galaxy clusters for decay, as shown in Figure~\ref{fig:comparison1}. This complementarity matters: pointed observations concentrate sensitivity on specific astrophysical targets selected by their brightness or J-factor, and their reach depends on assumptions about the DM profile of the targeted object. Cross-correlations integrate over the entire unresolved DM distribution and are insensitive to the DM profile of any single halo; they probe the aggregate signal from the large-scale structure. A detection in both channels would be strong evidence for a diffuse DM component tracing the cosmic web --- a qualitatively different signature from one localised in a single target. We now turn to the paper for the full derivation and results.
